\documentclass[11pt,a4paper]{article}
\usepackage{amsmath,amsfonts,amssymb,amsthm}
\usepackage{geometry}
\geometry{margin=1in}
\usepackage{graphicx}
\usepackage{hyperref}
\usepackage{xcolor}
\usepackage{enumitem}

\newtheorem{definition}{Definition}
\newtheorem{theorem}{Theorem}
\newtheorem{proposition}{Proposition}

\title{A Harmonic Sine Transform and the Riemann Zeta Function:\\
A Research Programme}
\author{}
\date{}

\begin{document}
\maketitle

\begin{abstract}
We outline a research programme that reformulates the classical Fourier transform in terms of the \emph{harmonic sine decomposition}—a greedy, trigonometric‑free algorithm that reconstructs $\sin(\theta)$ from a sequence of harmonic fractions $\pi/(n+2)$. The programme defines a \emph{harmonic sine transform} (HST) based on the greedy selection indicators $\delta_n(\theta)$ and their centred variants $e_n(\theta)$. The resulting non‑orthogonal basis possesses a natural arithmetic structure that links it directly to the Dirichlet series of the Riemann zeta function. The ultimate goal is to exploit this link by constructing a zeta‑weighted kernel whose spectral properties encode the distribution of the zeros of $\zeta(s)$, thereby providing a novel harmonic‑sine pathway toward the Riemann Hypothesis.
\end{abstract}

\section{Introduction}
The classical Fourier transform decomposes a function into orthogonal sinusoidal modes $\{ \sin(n\theta), \cos(n\theta) \}$. While the sine and cosine functions are analytic, their definition already presupposes the concept of angle. Victor Geere’s recent work~\cite{Geere2026a,Geere2026b} presents a radically different approach: a greedy algorithm that builds $\sin(\theta)$ purely from the Pythagorean theorem and a sequence of \emph{harmonic fractions}

\[
\alpha_n = \frac{\pi}{n+2}, \qquad n = 0,1,2,\dots
\]

The algorithm produces a family of step functions—the greedy selection indicators $\delta_n(\theta)$—that converge to $\sin(\theta)$ in the limit.

This research programme asks: \emph{can we replace the Fourier basis by the harmonic‑sine basis and study the resulting transform in relation to the Riemann zeta function?} The answer appears to be yes: the harmonic fractions $\pi/(n+2)$ are the terms of a scaled harmonic series, and the selection indicators $\delta_n(\theta)$ trace a path through the unit circle that is intimately connected with the logarithmic derivative of $\zeta(s)$.

\section{The Harmonic Sine Decomposition}
\subsection{Greedy harmonic angles}
For any angle $\theta \in [0,\pi]$ define the cumulative sum
\[
\theta_0 = 0, \qquad
\theta_{n+1} = \theta_n + \delta_n(\theta)\,\alpha_n,
\]
where the greedy selection indicator is
\[
\delta_n(\theta) = \Theta\bigl( \theta - \theta_n - \alpha_n \bigr) \in \{0,1\},
\]
and $\Theta$ is the Heaviside step function.
Geere’s main theorem states that the sequence
\[
s_n = \sin(\theta_n), \qquad n = 0,1,2,\dots
\]
converges to $\sin(\theta)$, with the error decaying as $O(1/n)$.

\subsection{The harmonic sine basis}
The indicators $\delta_n(\theta)$ are the fundamental building blocks. To obtain a zero‑mean basis suitable for harmonic analysis we centre them:
\[
e_n(\theta) = \delta_n(\theta) - \frac{\theta}{\pi}, \qquad \theta \in [0,\pi].
\]
Each $e_n$ is a piecewise‑constant function that jumps at a unique threshold angle $\theta_n^*$. The family $\{e_n\}_{n=0}^\infty$ is complete in $L^2[0,\pi]$ and possesses the correlation kernel
\[
K(n,m) = \int_0^\pi e_n(\theta)\, e_m(\theta)\, d\theta,
\]
which is positive semi‑definite but dense, reflecting the non‑orthogonal nature of the basis.

\section{The Harmonic Sine Transform}
\begin{definition}[Harmonic sine transform]
For $f \in L^2[0,\pi]$ the \emph{harmonic sine transform} (HST) is the sequence $(a_n)_{n=0}^\infty$ given by
\[
a_n = \langle f, e_n \rangle = \int_0^\pi f(\theta)\, e_n(\theta)\, d\theta.
\]
\end{definition}

The transform satisfies a Parseval‑type identity:
\[
\int_0^\pi f(\theta)\, g(\theta)\, d\theta
   = \sum_{n,m=0}^\infty a_n \, (K^{-1})_{nm} \, b_m,
\]
where $b_m = \langle g, e_m \rangle$ and $K^{-1}$ is the inverse of the Gram matrix. Inversion is therefore accomplished by the dual basis $\tilde{e}_n = \sum_m (K^{-1})_{nm}\, e_m$:
\[
f(\theta) = \sum_{n=0}^\infty a_n \, \tilde{e}_n(\theta).
\]

The HST differs from the classical Fourier sine transform in three essential ways:
\begin{enumerate}[label=(\roman*)]
\item \textbf{Adaptivity.} The basis functions are not periodic sinusoids but step functions whose jumps are determined by the greedy harmonic algorithm.
\item \textbf{Non‑orthogonality.} The Gram matrix is dense, requiring the use of the dual basis for reconstruction.
\item \textbf{Arithmetic sensitivity.} The harmonic fractions $\alpha_n = \pi/(n+2)$ embed the harmonic series directly into the basis, making the transform naturally attuned to number‑theoretic functions.
\end{enumerate}

\section{Connection to the Riemann Zeta Function}
\subsection{The harmonic series and $\zeta(s)$}
The harmonic fractions $\alpha_n = \pi/(n+2)$ are the terms of the scaled harmonic series. Consequently, for any fixed $\theta$, the Dirichlet weighted sum of the selection indicators,
\[
\sum_{n=0}^\infty \frac{\delta_n(\theta)}{(n+2)^s}, \qquad \operatorname{Re} s > 1,
\]
converges absolutely and defines a holomorphic function of $s$. When $\theta$ is integrated against the harmonic sine basis, we obtain a family of positive‑definite kernels
\[
K_s(n,m) = \int_0^\pi e_n(\theta)\, e_m(\theta)
           \left( \sum_{k=0}^\infty \frac{\delta_k(\theta)}{(k+2)^s} \right) d\theta.
\]
For $\operatorname{Re} s > 1$, $K_s$ is a holomorphic matrix‑valued function that inherits the analytic continuation of the Riemann zeta function through the identity
\[
\sum_{k=0}^\infty \frac{1}{(k+2)^s} = \zeta(s) - 1.
\]

\subsection{Relation to the explicit formula}
The logarithmic derivative of $\zeta(s)$ possesses the Dirichlet series
\[
-\frac{\zeta'}{\zeta}(s) = \sum_{n=1}^\infty \frac{\Lambda(n)}{n^s}, \qquad \operatorname{Re} s > 1,
\]
where $\Lambda(n)$ is the von Mangoldt function. Replacing the continuous Fourier modes in the Weil explicit formula by the harmonic sine basis leads to the correspondence
\[
\sum_{\rho} \widehat{f}(\rho) \longleftrightarrow
\sum_{n=0}^\infty \frac{a_n}{(n+2)^\rho},
\]
where $a_n$ are the HST coefficients of a suitably chosen test function $f$. The zeros $\rho$ of $\zeta(s)$ therefore appear as \emph{resonances} of the harmonic sine transform: they are the exponents for which the weighted sum of HST coefficients vanishes.

More concretely, for a test function $f$ whose HST is supported on a finite set of indices, the explicit formula becomes a relation between the zeros of $\zeta(s)$ and the prime‑indexed von Mangoldt weights, mediated entirely by the harmonic sine basis. This opens the prospect of isolating the contribution of individual zeros through a judicious choice of the test function.

\section{Outlook: Toward a Harmonic Sine Proof of the Riemann Hypothesis}
The research programme envisions the following computational and theoretical objectives:

\begin{enumerate}[start=1]
\item \textbf{Explicit construction of the dual basis.} Develop efficient numerical methods for computing the dual basis $\tilde{e}_n$ via the inverse Gram matrix, enabling fast harmonic sine transforms.

\item \textbf{Harmonic sine decomposition of $\zeta(1/2+it)$.} Compute the HST coefficients of the Riemann zeta function on the critical line. Preliminary experiments suggest that the coefficients exhibit a striking decay pattern that correlates with the zero ordinates $\gamma_k$.

\item \textbf{Zeta‑weighted kernel and the trace formula.} Study the kernel $K_s(n,m)$ as a function of $s$ and identify its spectral decomposition. The poles of $K_s$ at $s = \rho$ (the non‑trivial zeros) should be reflected in the null‑space of the kernel, providing a harmonic‑sine analogue of the Riemann–Weil explicit formula.

\item \textbf{Positivity criteria.} Investigate whether the positivity of the harmonic sine correlation kernel for $\operatorname{Re} s > 1/2$ is equivalent to the Riemann Hypothesis. This would generalise Weil’s positivity criterion to the harmonic sine setting.

\item \textbf{Connection to the Hilbert–Pólya conjecture.} Explore whether the harmonic sine basis can be used to construct a self‑adjoint operator whose eigenvalues are the imaginary parts of the zeta zeros, thereby realising the Hilbert–Pólya programme within a harmonic‑sine framework.
\end{enumerate}

\section*{Acknowledgments}
The author thanks Victor Geere for making the harmonic sine decomposition publicly available and for stimulating discussions on the harmonic‑sine reconstruction of trigonometric functions.

\begin{thebibliography}{9}
\bibitem{Geere2026a}
V.~Geere, \emph{A Purely Geometric Reconstruction of the Sine Function},
March~2026. \url{https://victorgeere.co.za/maths/sine-harmonic.html}
\bibitem{Geere2026b}
V.~Geere, \emph{On the Correlation Structure of a Greedy Harmonic Decomposition of the Sine Function},
March~2026. \url{https://victorgeere.co.za/maths/greedy-harmonic-decomposition.html}
\end{thebibliography}

\end{document}